% =====================================================================
%  TONG HOP THIET KE - USECASE: CHINH SUA YEU CAU NHAP HANG
% ---------------------------------------------------------------------
%  HUONG DAN BIEN DICH (QUAN TRONG):
%   - Bien dich bang XeLaTeX (KHONG dung pdfLaTeX vi co tieng Viet Unicode).
%   - Local (MiKTeX/TeX Live):   xelatex TongHop_ThietKe_YeuCauNhapHang.tex
%                                 (chay 2 lan de cap nhat Muc luc)
%   - Overleaf:  Menu -> Compiler -> chon "XeLaTeX".
%  Font "TeX Gyre Termes" (giong Times New Roman) co san glyph tieng Viet.
%  => Giu nguyen preamble nay de DONG BO style voi ban SRS goc.
% =====================================================================
\documentclass[12pt,a4paper]{article}

% ----- Font & Unicode (XeLaTeX) -----
\usepackage{fontspec}
\setmainfont{TeX Gyre Termes}      % Times-like, ho tro tieng Viet
\setsansfont{TeX Gyre Heros}

% ----- Trang & bo cuc -----
\usepackage{geometry}
\geometry{a4paper, margin=2.5cm}
\usepackage{amsmath,amssymb}
\usepackage{graphicx}
\usepackage{float}
\usepackage[table]{xcolor}
\usepackage{array}
\usepackage{tabularx}
\usepackage{longtable}
\usepackage{pdflscape}
\usepackage{enumitem}
\usepackage{listings}
\usepackage{fancyhdr}
\usepackage[hidelinks]{hyperref}
\hypersetup{unicode=true, colorlinks=true, linkcolor=blue!60!black,
            urlcolor=blue!60!black, citecolor=blue!60!black}

% ----- Ten nhan tieng Viet -----
\renewcommand{\contentsname}{Mục lục}
\renewcommand{\figurename}{Hình}
\renewcommand{\tablename}{Bảng}
\XeTeXlinebreaklocale "vi"
\XeTeXlinebreakskip = 0pt plus 0.1pt

% ----- Kieu cot bang -----
\newcolumntype{Y}{>{\raggedright\arraybackslash}X}
\newcolumntype{C}{>{\centering\arraybackslash}X}
\renewcommand{\arraystretch}{1.25}

% ----- Header / Footer -----
\pagestyle{fancy}
\fancyhf{}
\lhead{\small Tổng hợp thiết kế: Chỉnh sửa yêu cầu nhập hàng}
\rhead{\small Nhóm 4}
\cfoot{\thepage}
\renewcommand{\headrulewidth}{0.4pt}

% ----- Kieu code Java -----
\definecolor{codekw}{rgb}{0.13,0.13,0.7}
\definecolor{codecm}{rgb}{0.0,0.45,0.0}
\definecolor{codestr}{rgb}{0.6,0.1,0.1}
\definecolor{codebg}{rgb}{0.97,0.97,0.97}
\lstset{
  language=Java,
  basicstyle=\ttfamily\footnotesize,
  keywordstyle=\color{codekw}\bfseries,
  commentstyle=\color{codecm}\itshape,
  stringstyle=\color{codestr},
  numbers=left, numberstyle=\tiny\color{gray}, numbersep=8pt,
  backgroundcolor=\color{codebg},
  frame=single, rulecolor=\color{gray!50},
  breaklines=true, showstringspaces=false, tabsize=2,
  captionpos=b, xleftmargin=14pt, framexleftmargin=14pt
}

% ----- Thong tin sinh vien -----
\newcommand{\svName}{Đỗ Phú Hưng}
\newcommand{\svMSSV}{20235738}
\newcommand{\svGVHD}{Trịnh Tuấn Đạt}
\newcommand{\svHP}{IT4549}
\newcommand{\svNhom}{Nhóm 4}

\begin{document}

% =====================================================================
%  TRANG BIA
% =====================================================================
\begin{titlepage}
\centering
\thispagestyle{empty}

{\large \textbf{ĐẠI HỌC BÁCH KHOA HÀ NỘI}}\\[2pt]
{\large \textbf{TRƯỜNG CÔNG NGHỆ THÔNG TIN VÀ TRUYỀN THÔNG}}\\[6pt]
\rule{\textwidth}{0.8pt}

\vspace{1.2cm}

\vspace{0.8cm}
{\Huge \textbf{TỔNG HỢP THIẾT KẾ}}\\[10pt]
{\LARGE \textbf{HỆ THỐNG ĐẶT HÀNG NHẬP KHẨU}}\\[18pt]
{\Large Usecase: \textbf{Chỉnh sửa yêu cầu nhập hàng}}\\[6pt]
{\large (Áp dụng nguyên lý SOLID \& các mẫu thiết kế hướng đối tượng)}

\vspace{1.6cm}
\begin{flushleft}\large
\hspace{1.5cm}
\begin{tabular}{@{}ll@{}}
\textbf{Giảng viên hướng dẫn:} & \svGVHD \\[8pt]
\textbf{Sinh viên thực hiện:}  & \svName \\[8pt]
\textbf{MSSV:}                 & \svMSSV \\[8pt]
\textbf{Mã số lớp học phần:}   & \svHP \\[8pt]
\textbf{Nhóm:}                 & \svNhom \\
\end{tabular}
\end{flushleft}

\vfill
{\large Hà Nội, 2026}
\end{titlepage}

% =====================================================================
%  MUC LUC
% =====================================================================
\tableofcontents
\thispagestyle{empty}
\newpage
\setcounter{page}{1}

% =====================================================================
\section{Mục tiêu tái thiết kế}

Phần này tổng hợp quá trình \textbf{tái cấu trúc (refactoring)} mã nguồn của
chức năng \emph{Chỉnh sửa yêu cầu nhập hàng}, nhằm áp dụng các nguyên lý thiết
kế (SOLID) và các mẫu thiết kế hướng đối tượng đã học. Mục tiêu cụ thể:

\begin{itemize}[leftmargin=*]
  \item Tách bạch trách nhiệm giữa giao diện, điều khiển, nghiệp vụ và truy
        xuất dữ liệu để mã dễ đọc, dễ bảo trì.
  \item Cô lập phần nghiệp vụ cốt lõi (\texttt{updateRequest}) khỏi giao diện
        và cơ sở dữ liệu để có thể \textbf{kiểm thử đơn vị} mà không cần DB thật.
  \item Loại bỏ trùng lặp, tăng khả năng tái sử dụng các thành phần giao diện.
\end{itemize}

Phạm vi tổng hợp giới hạn trong gói (package):
\begin{center}
\texttt{com.app.modules.sales.request.editrequest}
\end{center}

% =====================================================================
\section{Hiện trạng trước khi tái thiết kế}

Ở phiên bản đầu, toàn bộ chức năng được dồn vào \textbf{một lớp giao diện duy
nhất} \texttt{EditRequestUI} (khoảng 440 dòng). Lớp này cùng lúc gánh nhiều
trách nhiệm:

\begin{itemize}[leftmargin=*]
  \item Hiển thị giao diện (render nhãn, bảng mặt hàng).
  \item Tự cấu hình từng cột và ô chỉnh sửa của bảng (\texttt{setupItemsTable},
        cùng các lớp lồng \texttt{QuantityCell}, \texttt{DeliveryDateCell}).
  \item Tự vẽ icon SVG bằng tay (\texttt{createIconButton}, \texttt{trashPath}).
  \item Tự bật các hộp thoại \texttt{Alert} xác nhận xóa, lưu, quay lại.
  \item Gọi thẳng tầng nghiệp vụ và chứa logic đồng bộ dữ liệu.
\end{itemize}

\noindent\textbf{Hệ quả:} vi phạm nghiêm trọng nguyên lý Đơn nhiệm (SRP) --- một
lớp có quá nhiều lý do để thay đổi; logic nghiệp vụ dính chặt với JavaFX nên
\emph{gần như không thể kiểm thử đơn vị}; các ô bảng viết riêng cho lớp này nên
không tái sử dụng được; mã icon SVG bị lặp lại.

% =====================================================================
\section{Kiến trúc sau khi tái thiết kế}

Một lớp khổng lồ được tách thành chuỗi lớp có trách nhiệm đơn nhất, tổ chức theo
mô hình \textbf{MVC} kết hợp \textbf{kiến trúc phân tầng} (Layered
Architecture). Luồng phụ thuộc một chiều: View $\rightarrow$ Controller
$\rightarrow$ Service $\rightarrow$ Repository $\rightarrow$ CSDL.

\begin{table}[H]
\centering
\caption{Phân rã trách nhiệm sau tái thiết kế}
\begin{tabularx}{\textwidth}{>{\ttfamily}l Y}
\hline
\rowcolor{gray!15}
\normalfont\textbf{Lớp / Thành phần} & \textbf{Trách nhiệm (một lý do thay đổi)} \\
\hline
EditRequestUI               & View thuần: giữ tham chiếu \texttt{@FXML}, hiển thị dữ liệu, phát sự kiện qua các hook. \\
EditRequestController       & Điều phối: bắt sự kiện, gọi Service, cập nhật View, gọi Alert. \\
EditRequestTableController  & Cấu hình cột và ô chỉnh sửa (cell factory) của bảng mặt hàng. \\
EditRequestAlertController  & Gom mọi hộp thoại cảnh báo / thông báo / xác nhận. \\
RequestService              & Nghiệp vụ: kiểm tra trạng thái \texttt{pending}, đồng bộ (upsert) mặt hàng. \\
RequestRepository           & Truy xuất CSDL: đọc/ghi bảng ImportRequest, RequestDetail. \\
UpdateRequestDTO, RequestResponse & DTO truyền dữ liệu qua ranh giới tầng. \\
IntegerEditCell, DatePickerCell, SvgButton & Thành phần giao diện dùng chung (\texttt{common/ui/components}). \\
\hline
\end{tabularx}
\end{table}

% =====================================================================
\section{Áp dụng nguyên lý SOLID}

\begin{longtable}{>{\bfseries}p{2.6cm} p{12.4cm}}
\hline
\rowcolor{gray!15}
\normalfont\textbf{Nguyên lý} & \textbf{Cách áp dụng trong mã nguồn} \\
\hline
\endhead
S --- Single Responsibility &
Mỗi lớp một lý do thay đổi: đổi giao diện sửa \texttt{EditRequestUI}; đổi cách
hiển thị bảng sửa \texttt{EditRequestTableController}; đổi quy tắc nghiệp vụ sửa
\texttt{RequestService}; đổi câu SQL sửa \texttt{RequestRepository}. \\
\hline
O --- Open/Closed &
\texttt{IntegerEditCell<S>}, \texttt{DatePickerCell<S>} là lớp generic nhận
hàm \texttt{getter}/\texttt{setter} qua lambda. Muốn dùng cho model khác chỉ cần
\emph{truyền tham số mới}, không sửa lớp ô --- mở để mở rộng, đóng với sửa đổi. \\
\hline
L --- Liskov Substitution &
Trong kiểm thử, \texttt{FakeRequestRepository extends RequestRepository} thay
thế hoàn toàn bản thật mà \texttt{RequestService} không cần biết, hành vi không
bị phá vỡ. \\
\hline
I --- Interface Segregation &
View không phụ thuộc một ``giao diện thần thánh''. Controller chỉ tương tác qua
các hook nhỏ: \texttt{Runnable saveAction}, \texttt{Consumer<RequestItem>
onDeleteItem} --- mỗi bên chỉ thấy đúng phần nó cần. \\
\hline
D --- Dependency Inversion &
\texttt{RequestRepository} phụ thuộc trừu tượng \texttt{IDBProvider} (không phụ
thuộc trực tiếp \texttt{PostgreSQLProvider}); \texttt{RequestService} nhận
\texttt{RequestRepository} qua constructor (Dependency Injection). \\
\hline
\end{longtable}

\noindent Minh họa nguyên lý Open/Closed và Dependency Injection:

\begin{lstlisting}[caption={Ô chỉnh sửa số nguyên dùng chung (generic, theo OCP)}]
public class IntegerEditCell<S> extends TableCell<S, S> {
    private final Function<S, Integer> getter;
    private final BiConsumer<S, Integer> setter;
    private final boolean editable;

    public IntegerEditCell(Function<S, Integer> getter,
                           BiConsumer<S, Integer> setter, boolean editable) {
        this.getter = getter; this.setter = setter; this.editable = editable;
        // ... lang nghe su kien commit gia tri
    }
}
\end{lstlisting}

\begin{lstlisting}[caption={Tiêm phụ thuộc qua constructor (DIP)}]
// Repository phu thuoc truu tuong IDBProvider, khong phu thuoc lop cu the
public RequestRepository() { this(new PostgreSQLProvider()); }
public RequestRepository(IDBProvider dbProvider) { this.dbProvider = dbProvider; }

// Service nhan Repository tu ben ngoai -> test bom duoc repo gia
public RequestService() { this(new RequestRepository()); }
public RequestService(RequestRepository repository) { this.repository = repository; }
\end{lstlisting}

% =====================================================================
\section{Các mẫu thiết kế đã áp dụng}

\begin{longtable}{>{\bfseries}p{3.0cm} p{4.0cm} p{7.8cm}}
\hline
\rowcolor{gray!15}
\normalfont\textbf{Mẫu thiết kế} & \textbf{Nơi áp dụng} & \textbf{Ý nghĩa} \\
\hline
\endhead
MVC & View / Controller / Model & Tách hiển thị -- điều khiển -- dữ liệu. \\
\hline
Layered (3 tầng) & UI $\to$ Service $\to$ Repository & UI không bao giờ chạm thẳng Repository. \\
\hline
DTO & \texttt{UpdateRequestDTO}, \texttt{RequestResponse} & Gói dữ liệu đi qua ranh giới tầng, tách entity nội bộ khỏi dữ liệu UI. \\
\hline
Repository & \texttt{RequestRepository} & Trừu tượng hóa truy xuất dữ liệu, ẩn SQL khỏi nghiệp vụ. \\
\hline
Dependency Injection & constructor \texttt{RequestService}, \texttt{RequestRepository} & Đảo ngược khởi tạo, phục vụ kiểm thử và mở rộng. \\
\hline
Factory + Singleton & \texttt{IDBProvider} / \texttt{PostgreSQLProvider} & Tập trung tạo kết nối, dùng chung một thể hiện. \\
\hline
Strategy (qua lambda) & getter/setter truyền vào \texttt{IntegerEditCell}, \texttt{DatePickerCell} & ``Cách lấy/ghi giá trị'' trở thành chiến lược thay thế được. \\
\hline
Command / Callback & \texttt{setActionForSaveButton(Runnable)}, \texttt{onSaved}, \texttt{onBack} & View phát sự kiện, không tự xử lý $\Rightarrow$ đảo phụ thuộc. \\
\hline
Observer & \texttt{ListChangeListener} cập nhật số mặt hàng & Giao diện tự đồng bộ khi dữ liệu đổi. \\
\hline
\end{longtable}

% =====================================================================
\section{So sánh trước / sau tái thiết kế}

\begin{table}[H]
\centering
\caption{Đối chiếu trạng thái mã nguồn}
\begin{tabularx}{\textwidth}{p{3.2cm} Y Y}
\hline
\rowcolor{gray!15}
\textbf{Tiêu chí} & \textbf{TRƯỚC} & \textbf{SAU} \\
\hline
Cấu trúc & 1 lớp \texttt{EditRequestUI} (~440 dòng) & View + 3 Controller + Service + Repository + DTO \\
\hline
Cấu hình ô bảng & lớp lồng viết cứng & \texttt{IntegerEditCell<S>}, \texttt{DatePickerCell<S>} generic, tái dùng \\
\hline
Hộp thoại & \texttt{Alert} rải khắp View & gom trong \texttt{EditRequestAlertController} \\
\hline
Truy xuất CSDL & trộn trong nghiệp vụ & tách \texttt{RequestRepository} + \texttt{IDBProvider} \\
\hline
Khả năng kiểm thử & gần như không & \texttt{RequestServiceTest} (Hộp đen + Hộp trắng C1) với repo giả \\
\hline
Khởi tạo phụ thuộc & \texttt{new} trực tiếp & Dependency Injection qua constructor \\
\hline
\end{tabularx}
\end{table}

% =====================================================================
\section{Ý nghĩa và lợi ích đem lại}

\begin{enumerate}[leftmargin=*]
  \item \textbf{Kiểm thử được (quan trọng nhất).} Nhờ DI và mẫu Repository, lớp
        \texttt{RequestServiceTest} kiểm thử trọn vẹn nghiệp vụ (upsert, chặn sửa
        khi \texttt{processing}/\texttt{completed}, xóa sạch mặt hàng) mà
        \emph{không cần CSDL thật}, dùng \texttt{FakeRequestRepository}. Bộ test
        phủ cả Hộp đen lẫn Hộp trắng (độ đo nhánh C1).
  \item \textbf{Dễ bảo trì.} Lỗi giao diện, lỗi nghiệp vụ và lỗi SQL nằm ở các
        lớp tách biệt; sửa một chỗ không gây vỡ chỗ khác.
  \item \textbf{Tái sử dụng.} \texttt{IntegerEditCell}, \texttt{DatePickerCell},
        \texttt{SvgButton} chuyển sang \texttt{common/ui/components} để các màn
        khác (tạo yêu cầu, xem chi tiết) dùng lại, xóa trùng lặp.
  \item \textbf{Dễ mở rộng.} Đổi nguồn dữ liệu chỉ cần một \texttt{IDBProvider}
        mới; thêm loại ô chỉnh sửa chỉ cần truyền lambda khác --- đúng tinh thần
        Open/Closed.
  \item \textbf{Giảm độ phức tạp nhận thức.} Từ một file dài ``làm tất cả''
        thành nhiều lớp nhỏ, mỗi lớp đọc hiểu trong vài phút.
\end{enumerate}

% =====================================================================
\section{Kết luận}

Việc tái thiết kế chức năng \emph{Chỉnh sửa yêu cầu nhập hàng} đã chuyển một lớp
giao diện đơn khối thành một kiến trúc phân tầng rõ ràng, tuân thủ đầy đủ năm
nguyên lý SOLID và vận dụng nhiều mẫu thiết kế kinh điển (MVC, Layered,
Repository, DTO, DI, Factory/Singleton, Strategy, Callback, Observer). Kết quả
là mã nguồn dễ đọc, dễ bảo trì, có khả năng tái sử dụng và --- quan trọng nhất
đối với mục tiêu kiểm thử --- cô lập được tầng nghiệp vụ để kiểm thử đơn vị đạt
độ phủ nhánh 100\% mà không phụ thuộc cơ sở dữ liệu.

\end{document}
